% ====================================================================
% EXPANDED SECTIONS 9, 10, 12, AND APPENDIX A
% For: Towards a Science for AI Coding Agent Harnesses
% ====================================================================
% This file contains expanded, self-contained versions of:
%   Section 9:  Cross-Pillar Synthesis
%   Section 10: Empirical Foundations
%   Section 12: Limitations and Open Problems
%   Appendix A: Cross-Pillar Theorem Index
% ====================================================================


% ====================================================================
% 9. CROSS-PILLAR SYNTHESIS
% ====================================================================
\section{Cross-Pillar Synthesis}
\label{sec:synthesis}

The preceding sections presented each pillar independently. This section, the primary contribution of this synthesis paper, identifies the structural connections that bind them into a unified framework.

\subsection{Information as the Unifying Currency}

The deepest structural connection across all seven pillars is that each reasons about \emph{information flow} through the harness:

\begin{itemize}[nosep]
    \item \textbf{Abstraction}: the information gap $K(P \mid S)$ between specification and code.
    \item \textbf{Information}: the Shannon decomposition $H(P \mid S) = H(P \mid S,C,\theta) + I(P;C \mid S,\theta) + I(P;\theta \mid S)$.
    \item \textbf{Reliability}: compound error as information destruction, $R(n) = p^n$.
    \item \textbf{Coordination}: error amplification as correlated information loss across agents.
    \item \textbf{Temporal}: context fidelity $\Phi(t) = e^{-\Lambda t}$ as information decay over time.
    \item \textbf{Quality}: the Detection Ceiling $d_{\ell,j} \leq I(X; D_j)/H(D_j)$ as an information bound on verification.
    \item \textbf{Governance}: governance capacity $C_{\text{gov}}$ as information throughput of the correction channel.
\end{itemize}

The unification is strongest along the \emph{semantic axis} (Abstraction, Information) and the \emph{assurance axis} (Quality, Governance), where information theory is the native formalism. For the \emph{execution axis} (Reliability, Coordination, Temporal), the information-theoretic framing is a valid interpretation but not the generative formalism: $R(n) = p^n$ is a probability product from reliability theory, the Extended Amdahl's Law draws on parallel computing and concurrency control, and VIH draws on scheduling and queueing theory. Calling compound errors ``information destruction'' is accurate (errors do corrupt the information content of pipeline output) but retrospective; the results would stand unchanged without the information-theoretic gloss. We adopt the information lens because it provides a common vocabulary for cross-pillar reasoning, while acknowledging that its explanatory power varies across pillars.

\subsection{The Abstraction Gap Chain}

The first structural connection links three pillars in a chain:

\begin{enumerate}
    \item \textbf{Abstraction} defines the gap: $\mathcal{G}(S,P) = K(P \mid S)$.
    \item \textbf{Information} operationalizes it: $H(P \mid S) = H(P \mid S,C,\theta) + I(P;C \mid S,\theta) + I(P;\theta \mid S)$, decomposing the gap into three actionable levers.
    \item \textbf{Quality} addresses what happens when the gap is traversed poorly: the 18 slop types are failure modes of the translation from spec to code, and the Compound Degradation Theorem (Theorem~\ref{thm:compound-degradation}) quantifies the consequences.
\end{enumerate}

The chain makes precise how a specification becomes code and what goes wrong along the way. The Abstraction pillar establishes the theoretical limits. The Information pillar identifies the levers. The Quality pillar classifies the failure modes and designs defenses.

\subsection{The Compound Error Cascade}

The second structural connection links three pillars through the theme of multiplicative degradation:

\begin{enumerate}
    \item \textbf{Reliability} proves $R(n) = p^n$: compound errors in sequential pipelines.
    \item \textbf{Coordination} extends to multi-agent settings: the Extended Amdahl's Law (Theorem~\ref{thm:extended-amdahl}) adds the reliability factor $(1-\delta_a)^n$, and empirical amplification ($17.2\times$) exceeds the independence prediction.
    \item \textbf{Temporal} frames the rate at which compound errors occur: VIH $= \lambda_{\text{raw}} \cdot p_{\text{verify}}$ captures the tradeoff between iteration speed and error rate.
\end{enumerate}

The cascade reveals a fundamental tension: the same exponential that makes compound errors devastating also makes per-step improvements highly leveraged (elasticity $= n$). This is the core design insight: invest in per-step quality, not pipeline sophistication.

\subsection{The Structural Enforcement Principle}

The principle that structural enforcement dominates instructional enforcement appears independently in four pillars:

\begin{table}[h]
\centering
\caption{The structural enforcement principle across pillars.}
\label{tab:structural}
\begin{tabular}{lp{10cm}}
\toprule
\textbf{Pillar} & \textbf{Manifestation} \\
\midrule
Information & Tier boundaries enforced structurally (file system, .gitignore) vs. instructionally (prompt rules). Compliance gap: $(1-\epsilon)^T$. \\
Reliability & Cost-of-failure threshold $L^*$ determines when structural enforcement is cost-effective. Threshold drops as $1/n$ for multi-step pipelines. \\
Quality & Decidable slop types should be caught by structural gates (Layer 1), not by probabilistic methods. \\
Coordination & File-ownership contracts enforce $F_i \cap F_j = \emptyset$ structurally, eliminating textual/syntactic conflicts by construction. \\
\bottomrule
\end{tabular}
\end{table}

The convergence is not coincidental. Structural enforcement converts probabilistic guarantees into deterministic ones. In any system where errors compound (as they do in all agent pipelines), converting even a small fraction of probabilistic steps to deterministic ones has outsized impact on end-to-end reliability.

\subsection{The Four-Layer Defense Architecture}

The four-layer verification architecture appears in both the Reliability and Quality pillars:

\begin{table}[h]
\centering
\caption{The shared four-layer verification architecture.}
\label{tab:four-layer}
\begin{tabular}{clccc}
\toprule
\textbf{Layer} & \textbf{Mechanism} & \textbf{Cost} & \textbf{Determinism} & \textbf{Decidability Class} \\
\midrule
1 & Structural gates & Very low & Deterministic & Decidable \\
2 & Automated testing & Low & Probabilistic & Semi-decidable \\
3 & AI review & Moderate & Probabilistic & Semi-decidable \\
4 & Human review & High & Judgment & Undecidable \\
\bottomrule
\end{tabular}
\end{table}

The ordering by cost and determinism is not arbitrary; it follows from the decidability classification. Decidable defects should be caught at Layer 1 (structural gates), where detection is certain and marginal cost is near zero. Semi-decidable defects require Layers 2--3, where detection is probabilistic but cost-effective. Undecidable defects require Layer 4, where human judgment provides the irreplaceable element.

\subsection{Governance Capacity as the Macro Constraint}

The Governance pillar's capacity bound (Theorem~\ref{thm:governance-capacity}) is the system-level analog of the Reliability pillar's compound error sensitivity (Theorem~\ref{thm:compound-sensitivity}):

\begin{itemize}[nosep]
    \item \textbf{Reliability}: at the step level, if $p < 1$, errors compound as $p^n$.
    \item \textbf{Governance}: at the system level, if $R_{\text{drift}} > C_{\text{gov}}$, divergence accumulates without bound.
\end{itemize}

Both are threshold phenomena: below the threshold, the system is self-correcting; above it, degradation is unbounded. Both have the same structural form: a rate of degradation versus a capacity for correction. And both yield the same design prescription: invest in increasing the correction capacity (per-step quality, governance bandwidth) rather than adding post-hoc repair mechanisms.

\subsection{Shared Formal Machinery}
\label{sec:shared-machinery}

Table~\ref{tab:machinery} maps the 16 distinct formal techniques used across the seven pillars. The density of the mapping (most techniques appear in 2+ pillars) reflects the deep structural unity of the framework; the pillars are not independent theories glued together but facets of a common problem viewed through different mathematical lenses.

\begin{table}[h]
\centering
\caption{Formal machinery shared across pillars. A bullet (\textbullet) indicates that the pillar makes substantive use of the technique; a circle ($\circ$) indicates secondary or analogical use.}
\label{tab:machinery}
\small
\begin{tabular}{l@{\hskip 4pt}c@{\hskip 4pt}c@{\hskip 4pt}c@{\hskip 4pt}c@{\hskip 4pt}c@{\hskip 4pt}c@{\hskip 4pt}c}
\toprule
\textbf{Technique} & \textbf{Abs} & \textbf{Inf} & \textbf{Rel} & \textbf{Crd} & \textbf{Tmp} & \textbf{Qlt} & \textbf{Gov} \\
\midrule
Shannon information theory ($H$, $I$, DPI) & \textbullet & \textbullet & $\circ$ & $\circ$ & $\circ$ & \textbullet & \textbullet \\
Kolmogorov complexity ($K$) & \textbullet & $\circ$ & & & & & \textbullet \\
Series reliability / compound errors ($p^n$) & & & \textbullet & \textbullet & $\circ$ & & \\
Structural vs. instructional enforcement & & \textbullet & \textbullet & \textbullet & & \textbullet & $\circ$ \\
Submodular optimization / greedy approx. & & \textbullet & & & & \textbullet & \\
Control theory (cascade, Nyquist, stability) & & & & & $\circ$ & & \textbullet \\
Lattice theory (Galois, closure, refinement) & \textbullet & & & & & & \textbullet \\
Concurrency control / isolation levels & & & & \textbullet & & & \\
Survival analysis (competing risks, Weibull) & & & & & & & \textbullet \\
Scheduling theory (Amdahl, stochastic DAG) & & & & \textbullet & \textbullet & & \\
FKG / correlation inequalities & & & \textbullet & & & \textbullet & \\
Graph partitioning (min-cut, Cheeger) & & & & \textbullet & & & \\
Rate-distortion theory & & & & & & & \textbullet \\
Queueing theory (Little's law, Kingman) & & & & & \textbullet & & \\
Speculative execution theory & & & & & \textbullet & & \\
Cache competitive analysis (Sleator-Tarjan) & & & & & \textbullet & & \\
\bottomrule
\end{tabular}
\end{table}

Three patterns emerge. First, Shannon information theory provides the broadest cross-pillar vocabulary, appearing substantively in four pillars and secondarily in three more. Second, structural enforcement is the most widely shared \emph{design principle} (four pillars derive it independently). Third, several techniques (concurrency control, survival analysis, cache competitive analysis) are used in only one pillar, suggesting that further cross-pollination may be productive; for instance, survival analysis could model the lifetime of context relevance (connecting Governance to Temporal), and concurrency control formalisms could sharpen the governance ratchet's interaction with parallel agent modifications.

\subsection{Shared Empirical Sources}
\label{sec:shared-empirical}

Table~\ref{tab:empirical-sources} maps the 14 primary empirical sources to the pillars that rely on them. The density of cross-pillar citations highlights both a strength (convergent evidence from multiple analytical perspectives) and a vulnerability (the same data points bear disproportionate evidentiary weight).

\begin{table}[h]
\centering
\caption{Key empirical sources and the pillars that draw on them. Columns indicate pillar usage: Abs(traction), Inf(ormation), Rel(iability), Crd (Coordination), Tmp (Temporal), Qlt (Quality), Gov(ernance).}
\label{tab:empirical-sources}
\small
\begin{tabular}{l@{\hskip 3pt}c@{\hskip 3pt}c@{\hskip 3pt}c@{\hskip 3pt}c@{\hskip 3pt}c@{\hskip 3pt}c@{\hskip 3pt}c@{\hskip 6pt}p{4.5cm}}
\toprule
\textbf{Source} & \textbf{Abs} & \textbf{Inf} & \textbf{Rel} & \textbf{Crd} & \textbf{Tmp} & \textbf{Qlt} & \textbf{Gov} & \textbf{Key Finding} \\
\midrule
GitClear 2025 & & \textbullet & & & \textbullet & \textbullet & \textbullet & 211M lines; 8$\times$ duplication; 60\% refactoring collapse \\
DORA 2025 & & & & \textbullet & \textbullet & \textbullet & \textbullet & Throughput up 21\%; stability negative \\
Kim et al.\ 2025 & \textbullet & & & \textbullet & & & & $17.2\times$ error amplification; 180 configs \\
METR 2026 & \textbullet & & \textbullet & & & \textbullet & & 50\% test-passing PRs unmerged \\
Liu et al.\ 2024 & & \textbullet & & & & & & Lost in the Middle; U-shaped recall \\
Chroma 2025 & & \textbullet & & & \textbullet & & & 18 models; monotonic degradation \\
JetBrains 2025 & & \textbullet & & & & & & Complexity increase in AI code \\
CodeRabbit 2025 & & & & & & \textbullet & & 1.68$\times$ issues; 2.74$\times$ XSS vulns \\
Veracode 2025 & & & & & & \textbullet & & AI code security defect rates \\
Apollo Research 2025 & & & \textbullet & & & \textbullet & & Scheming behaviors in agents \\
NIST CAISI 2025 & & & \textbullet & & & \textbullet & & Benchmark cheating evidence \\
Hariharan et al.\ 2025 & & & \textbullet & & & & & 63\% failure on 100-step tasks \\
Spotify Honk 2025 & & \textbullet & & & & \textbullet & & Production harness telemetry \\
Gloaguen et al.\ 2026 & & \textbullet & & & & & & AGENTS.md convention data \\
\bottomrule
\end{tabular}
\end{table}

The table reveals a concentration risk: GitClear, DORA, and METR each support claims in 3--4 pillars, meaning that a successful challenge to any one of these sources would ripple across the framework. Section~\ref{sec:empirical} provides a quality assessment of each source to make this risk explicit.

\subsection{Additional Cross-Pillar Connections}
\label{sec:missed-connections}

Beyond the three primary cascades (abstraction gap chain, compound error cascade, structural enforcement principle), five cross-pillar interactions merit formal development.

\paragraph{1. Accretion meets the ratchet.}
The Quality pillar's Accretion Category (individually CI-passing changes that collectively degrade codebase health) is precisely the failure mode that the Governance pillar's ratchet mechanism (Definition~\ref{def:ratchet}) aims to prevent. The interaction, however, is non-trivial. Accretion defects are correct by construction; they satisfy all existing ratchet constraints. To prevent accretion, the ratchet must enforce \emph{structural health metrics} (duplication rate, dependency fan-out, abstraction depth) rather than per-change correctness. But structural health constraints are inherently more brittle than correctness constraints: they have higher false-positive rates and are sensitive to legitimate refactoring. Adding ratchet constraints on structural metrics risks triggering the ossification pathology (the Governance pillar's Ossification Theorem), where the constraint set grows until the allowed change set becomes empty or effectively so. The calibration question is: at what threshold should structural health metrics be ratcheted, and how should evidence expiry windows be set relative to the accretion rate $\alpha(t)$? The Compound Degradation Theorem (Theorem~\ref{thm:compound-degradation}) provides the growth rate; the Ratchet Convergence theorem provides the stabilization guarantee; what is missing is the bridging analysis showing how to set ratchet parameters that prevent accretion without inducing ossification.

\paragraph{2. Staleness compounds with coordination.}
The Temporal pillar's context fidelity decay $\Phi(t) = e^{-\Lambda(t-t_0)}$ interacts multiplicatively with the Coordination pillar's error amplification. In multi-agent settings, the effective decay rate $\Lambda$ is not a constant; it increases with the number of concurrent agents, because more agents means faster codebase evolution and therefore faster context invalidation. Specifically, if $k$ agents each commit at rate $\lambda_c$, the effective decay rate for any one agent's context is $\Lambda_{\text{eff}} \approx (k-1)\lambda_c \alpha$, where $\alpha$ is the fraction of commits that modify files in the agent's active context. This creates a feedback loop: more agents $\to$ faster staleness $\to$ higher per-agent error rate $\to$ more rework $\to$ more commits $\to$ even faster staleness. The Extended Amdahl's Law (Theorem~\ref{thm:extended-amdahl}) does not account for this feedback; incorporating staleness-induced error rate inflation into the reliability factor $(1-\delta_a)^n$ would likely reduce $n^*$ below the current estimate $1/\sqrt{\delta_a}$, particularly for tightly coupled codebases where $\alpha$ is large.

\paragraph{3. Lattice parallels across Abstraction and Governance.}
Both the Abstraction pillar and the Governance pillar employ lattice-theoretic structures, but with different orientations. The Abstraction pillar's refinement lattice orders specifications from most abstract (top, natural language) to most concrete (bottom, executable code), with the Galois connection $(\alpha, \gamma)$ mediating between specifications and programs. The Governance pillar's constraint lattice orders constraint sets by inclusion, with the ratchet as a closure operator that moves the system toward the join (more constrained). These two lattices are not independent: every specification in the refinement lattice \emph{induces} a constraint in the governance lattice (the specification constrains which programs are acceptable). A formal bridge would establish that the Galois connection on the refinement lattice is compatible with the closure operator on the constraint lattice; that is, refining a specification (descending the refinement lattice) should tighten constraints (ascending the constraint lattice) monotonically. If this compatibility holds, then the act of specification refinement automatically generates governance constraints, providing a formal foundation for ``specification as governance.''

\paragraph{4. Decomposition enables reuse discovery.}
The Coordination pillar's task decomposition (balanced min-cut on the coupling graph) and the Information pillar's reuse discovery problem are connected through the structure of the coupling graph itself. Balanced min-cut partitions identify clusters of highly coupled files; these clusters are also the natural units for reuse discovery, because semantic similarity concentrates within clusters. An agent assigned to a cluster via the decomposition has, by construction, the most relevant context for discovering reuse opportunities within that cluster. This suggests that coordination topology should be co-designed with the context architecture: the partition that minimizes merge conflicts (Coordination) simultaneously maximizes reuse-relevant context concentration (Information). The formal connection is that the coupling graph's cluster structure determines both the optimal decomposition and the natural reuse neighborhoods, and a single spectral analysis of the coupling graph can serve both objectives.

\paragraph{5. FKG amplification across coordination boundaries.}
The Quality pillar's FKG inequality (Proposition~\ref{prop:fkg}), which establishes that same-family LLM judges underestimate escape probabilities due to correlated blind spots, has a multi-agent analog that the current framework does not address. When $k$ agents from the same model family work on a shared codebase, their errors are positively correlated not only within each agent's pipeline (the Quality pillar's concern) but \emph{across} agents (the Coordination pillar's concern). The FKG inequality implies:
\begin{equation}
    P\!\left(\bigcap_{i=1}^k E_i\right) \geq \prod_{i=1}^k P(E_i)
\end{equation}
where $E_i$ is the event that agent $i$ introduces a defect in the shared semantic category. This means the Coordination pillar's independence assumption ($A_e \to k$) not only underestimates amplification (as the Kim et al.\ empirical data shows) but does so in a structurally predictable way: the FKG framework provides an \emph{upper} bound on the independence gap. Combining the FKG inequality with the coordination error amplification model would yield tighter bounds on multi-agent defect rates, explaining part of the $6\times$ discrepancy between the theoretical independence prediction and the empirical $17.2\times$ amplification.

\subsection{The Governance Information Budget}
\label{sec:gib}

We propose the \emph{Governance Information Budget} as the integrating concept across all seven pillars. The idea is that every harness operates under a total information budget, and the pillars decompose how that budget is allocated:

\begin{equation}
\label{eq:gib}
    \underbrace{K(P \mid S)}_{\text{Abstraction gap}} = \underbrace{I(P; C \mid S, \theta)}_{\text{Context (Information)}} + \underbrace{I(P; \theta \mid S)}_{\text{Prior (Model)}} + \underbrace{H(P \mid S, C, \theta)}_{\text{Residual (Harness must handle)}}
\end{equation}

The residual term $H(P \mid S, C, \theta)$ is the information that neither the context nor the model provides. This residual must be filled by the agent's exploration and reasoning, which is subject to:
\begin{itemize}[nosep]
    \item Compound error sensitivity (Reliability): each exploratory step has failure probability $q$, compounding as $(1-q)^n$.
    \item Error amplification (Coordination): multi-agent exploration amplifies errors by up to $17.2\times$.
    \item Context staleness (Temporal): the residual grows as context decays, $\Phi(t) = e^{-\Lambda t}$.
    \item Detection limits (Quality): some errors in filling the residual are undetectable, $d_{\ell,j} \leq I(X; D_j)/H(D_j)$.
    \item Governance capacity (Governance): the total rate at which the residual is filled must not exceed governance capacity $C_{\text{gov}}$.
\end{itemize}

\begin{proposition}[Governance Information Budget Constraint]
\label{prop:gib-constraint}
For a harness operating at steady state with change rate $\lambda$ (changes per unit time), the Governance Information Budget imposes five simultaneous constraints, one from each execution-axis and assurance-axis pillar:
\begin{align}
    \text{(Reliability)} &\quad \lambda \leq \frac{-\ln R_{\min}}{n \cdot |\ln p|} \label{eq:gib-rel} \\
    \text{(Coordination)} &\quad k \leq n^* \approx 1/\sqrt{\delta_a} \label{eq:gib-coord} \\
    \text{(Temporal)} &\quad \lambda \leq \Lambda / \ln(1/\Phi_{\min}) \label{eq:gib-temp} \\
    \text{(Quality)} &\quad \lambda \cdot P_{\text{esc}} \leq \mu_{\text{repair}} \label{eq:gib-qual} \\
    \text{(Governance)} &\quad \lambda \cdot H_{\text{drift/change}} \leq C_{\text{gov}} \label{eq:gib-gov}
\end{align}
where $R_{\min}$ is the minimum acceptable pipeline reliability, $\Phi_{\min}$ is the minimum acceptable context fidelity, $P_{\text{esc}}$ is the escape probability past all defense layers, $\mu_{\text{repair}}$ is the defect repair rate, and $H_{\text{drift/change}}$ is the architectural drift per change (in bits). The binding constraint (the one that limits $\lambda$ most severely) determines which pillar is the system's current bottleneck.
\end{proposition}

\begin{remark}
The five constraints in Proposition~\ref{prop:gib-constraint} are not independent. Increasing $k$ (more parallel agents) relaxes the Reliability constraint (more throughput) but tightens the Temporal constraint (faster staleness) and the Governance constraint (more drift per unit time). Increasing per-step quality $p$ relaxes the Reliability constraint but has no direct effect on the Governance constraint. This interdependence is the core reason that harness design requires simultaneous optimization across pillars, not sequential optimization of each pillar independently.
\end{remark}

The Governance Information Budget unifies the pillars by showing that they all constrain different aspects of the same fundamental problem: closing the abstraction gap reliably, efficiently, and coherently. The binding constraint rotates as the harness configuration changes: early in development (low $k$, high residual), the Abstraction and Information pillars dominate; at scale (high $k$, fast $\lambda$), the Coordination and Governance pillars dominate; for mature, stable systems (low $\lambda$, long-lived context), the Temporal pillar dominates through staleness accumulation.


% ====================================================================
% 10. EMPIRICAL FOUNDATIONS
% ====================================================================
\section{Empirical Foundations}
\label{sec:empirical}

The formal framework developed in Sections~\ref{sec:abstraction}--\ref{sec:synthesis} rests on empirical evidence of varying quality. An honest assessment of this evidence is essential.

\subsection{Strong Empirical Foundations}

\paragraph{GitClear 2025.} Analysis of 211 million changed lines across a multi-year period \citep{gitclear2025}. Key findings: code blocks with 5+ duplicated lines increased approximately 8$\times$ (overall duplication rate rose from 8.3\% to 12.3\%); refactoring collapsed from roughly 25\% to under 10\% of changed lines. Strengths: large sample, longitudinal design, objective metrics. Limitations: observational (no causal identification), single platform, code-level metrics only (no quality assessment).

\paragraph{DORA 2025.} Telemetry from over 10,000 developers \citep{dora2025}. Key finding: individual task completion increased 21\% with AI adoption, but organizational-level throughput stayed flat and software delivery stability showed a negative correlation with AI usage. Strengths: large sample, standardized metrics, multi-organization. Limitations: survey-based components, correlation not causation, no controlled intervention.

\paragraph{Kim et al.\ 2025.} 180 configurations across 5 architectures, 3 LLM families, and 4 benchmarks \citep{kim2025scaling}. Key finding: $17.2\times$ error amplification for independent agents, saturation of scaling benefits at approximately 45\% of the task completion rate. Strengths: systematic, controlled, comprehensive coverage of configuration space. Limitations: benchmark tasks (not production code), specific models and benchmarks.

\paragraph{METR 2026.} Expert evaluation of SWE-bench test-passing pull requests \citep{metr2026swebench}. Key finding: roughly 50\% of test-passing PRs would not be merged by repository maintainers (but only 68\% of human-written golden patches would be merged on re-review, so the AI-attributable gap is approximately 18 percentage points). Strengths: expert evaluation, direct relevance to the circular validation problem. Limitations: subjective judgments, small evaluator pool (4 maintainers, 3 repositories), specific benchmark context.

\subsection{Moderate Empirical Foundations}

\paragraph{Context degradation studies.} Multiple studies \citep{chroma2025contextrot, liu2024lost, du2025context} demonstrating monotonic performance degradation with context length. The power-law model $\delta(n) = (n/W_{\text{eff}})^{\alpha_d}$ with $\alpha_d \in [0.3, 0.5]$ is a fit to this data, not a derived result. The fit is reasonable but not definitive.

\paragraph{CodeRabbit 2025.} Analysis of 470 pull requests \citep{coderabbit2025}. Findings: 1.68$\times$ more issues in AI-generated PRs; 2.74$\times$ more XSS vulnerabilities. Strengths: direct comparison, automated analysis. Limitations: modest sample size, single tool, potential selection bias.

\paragraph{Hariharan et al.\ 2025.} Evaluation of plan-and-execute agent architectures \citep{hariharan2025plan}. Key finding: 63\% failure rate on 100-step tasks. Strengths: systematic scaling study. Limitations: specific agent architecture, benchmark-only evaluation.

\subsection{Weak or Missing Empirical Foundations}

Several key claims in the framework lack strong empirical support:

\begin{itemize}
    \item The optimal context budget ($|C^*| \approx 0.15W$ to $0.40W$) is derived from model fits, not from direct optimization experiments.
    \item The structural enforcement compliance gap $(1-\epsilon)^T$ is a theoretical prediction; no controlled study has measured $\epsilon$ for different enforcement mechanisms.
    \item The governance capacity bound (Theorem~\ref{thm:governance-capacity}) is a theoretical result by construction; calibrating $R_{\text{drift}}$ and $C_{\text{gov}}$ for real systems remains an open challenge.
    \item The Accretion Category is well-motivated by observational data but has not been validated through controlled experiments.
    \item VIH has not been measured in production harness settings.
    \item The cache-staleness EOQ theorem has not been empirically validated.
\end{itemize}

\subsection{Evidence Quality Assessment}
\label{sec:evidence-quality}

Table~\ref{tab:evidence-quality} provides a systematic quality assessment of the empirical foundations supporting the framework. Quality ratings follow a four-tier scheme: \textbf{High} (large sample, controlled or quasi-experimental, replicable); \textbf{Medium} (moderate sample, systematic but observational, partially replicable); \textbf{Low} (small sample, anecdotal or single-platform, limited replicability); \textbf{Very Low} (theoretical prediction only, no direct empirical measurement).

\begin{table}[H]
\centering
\caption{Evidence quality assessment for major empirical sources and theoretical claims. ``Key Gap'' identifies the most important missing piece for each data category.}
\label{tab:evidence-quality}
\small
\begin{tabular}{p{3cm}cccp{4.2cm}}
\toprule
\textbf{Data Category} & \textbf{Quality} & \textbf{Sample} & \textbf{Design} & \textbf{Key Gap} \\
\midrule
GitClear (duplication, refactoring) & High & 211M lines & Longitudinal, observational & No causal identification; single platform; no direct quality measure \\
\addlinespace
DORA 2025 (throughput, stability) & High & 10,000+ devs & Cross-sectional survey + telemetry & Survey components; no controlled intervention; correlation only \\
\addlinespace
Kim et al.\ (error amplification) & High & 180 configs & Controlled, systematic & Benchmark tasks only; 3 model families; not production code \\
\addlinespace
METR (merge gap) & Medium & $\sim$100 PRs & Expert evaluation & 4 evaluators, 3 repos; subjective criteria; small sample \\
\addlinespace
Liu et al.\ (Lost in Middle) & Medium & Multiple models & Controlled, synthetic tasks & Synthetic needle tasks; unclear production generalization \\
\addlinespace
Chroma (context rot) & Medium & 18 models & Controlled, multi-provider & Synthetic benchmarks; degradation model is a fit, not derived \\
\addlinespace
JetBrains (complexity) & Medium & Survey + metrics & Cross-sectional & Self-reported AI usage; single IDE ecosystem \\
\addlinespace
CodeRabbit (PR quality) & Medium & 470 PRs & Comparative & Single tool; potential selection bias; modest sample \\
\addlinespace
Veracode (security) & Medium & Undisclosed & Static analysis comparison & Methodology details limited; no controlled experiment \\
\addlinespace
Apollo Research (scheming) & Low & Case studies & Behavioral evaluation & Small-scale; specific models; adversarial setup \\
\addlinespace
NIST CAISI (cheating) & Low & Workshop report & Expert panel consensus & No quantitative data; policy-oriented rather than empirical \\
\addlinespace
Optimal context budget & Very Low & None & Model fit to degradation data & No direct optimization experiment; parameter sensitivity unknown \\
\addlinespace
Structural enforcement $\epsilon$ & Very Low & None & Theoretical prediction & No measurement of per-step violation probability \\
\addlinespace
VIH in production & Very Low & None & Proposed metric, unmeasured & No harness instrumentation; no baseline data \\
\addlinespace
Governance capacity ($C_{\text{gov}}$) & Very Low & None & Theoretical construction & $R_{\text{drift}}$ and $C_{\text{gov}}$ undefined operationally \\
\addlinespace
Accretion detection & Very Low & None & Category proposed, unvalidated & No labeled dataset; no precision/recall measurement \\
\addlinespace
Cache-staleness EOQ & Very Low & None & Theoretical analog & No measurement of staleness cost or optimal refresh interval \\
\bottomrule
\end{tabular}
\end{table}

\subsection{The Calibration Gap}

A recurring pattern across all seven pillars is the gap between theoretical models and empirical calibration. The models are well-grounded in established formal machinery, and the qualitative predictions align with practitioner experience. But the quantitative parameters (decay rates, threshold values, optimal counts) are either estimated from limited data or derived from model assumptions rather than measured directly.

This gap is not a criticism of the framework; it is a statement of the current state of the field. Closing this gap requires controlled experiments, which in turn requires access to production harness telemetry. Section~\ref{sec:limitations} proposes specific experiments, organized as a concrete verification roadmap with falsification criteria for each major formal result.


% ====================================================================
% 12. LIMITATIONS AND OPEN PROBLEMS
% ====================================================================
\section{Limitations and Open Problems}
\label{sec:limitations}

\subsection{Empirical Validation Gaps}

The most significant limitation of this framework is the gap between formal models and empirical validation. The following experiments are needed to close this gap:

\begin{enumerate}
    \item \textbf{Context budget validation}: controlled experiments measuring task success as a function of context size, across multiple models and task types, to validate the optimal budget estimates ($|C^*| \approx 0.15W$ to $0.40W$).
    \item \textbf{Structural enforcement measurement}: controlled comparison of instructional vs. structural enforcement, measuring compliance rates $\epsilon$ across different enforcement mechanisms and pipeline lengths $T$.
    \item \textbf{VIH measurement}: instrumentation of production harnesses to measure raw iteration rate, verification pass rate, and their product across different verification configurations.
    \item \textbf{Governance capacity calibration}: measurement of $R_{\text{drift}}$ and $C_{\text{gov}}$ in production settings, to test whether the capacity bound predicts the onset of architectural degradation.
    \item \textbf{Accretion detection}: development and validation of detectors for the Accretion Category, measuring precision and recall against expert-labeled datasets.
\end{enumerate}

\subsection{Empirical Verification Roadmap}
\label{sec:verification-roadmap}

This subsection provides a concrete verification roadmap for the five most important formal results in the framework. For each result, we specify the data to collect, the testable prediction, and the falsification criterion. This roadmap was developed in direct response to peer review feedback identifying the empirical gap as the framework's most significant weakness.

\subsubsection{Result 1: Compound Error Sensitivity (Theorem~\ref{thm:compound-sensitivity})}

\paragraph{Formal claim.} For an $n$-step agent pipeline with homogeneous per-step success probability $p$, the end-to-end reliability is $R(n) = p^n$, and the elasticity of $R$ with respect to $p$ equals the pipeline length: $\mathcal{E}(R,p) = n$.

\paragraph{Data to collect.}
\begin{enumerate}[nosep]
    \item Instrument a production agent harness (e.g., Claude Code, Cursor, Codex) to log, for each pipeline execution: (a) the number of discrete steps $n$, (b) a binary success/failure indicator for each step (judged by automated tests plus human spot-checks on a stratified subsample), and (c) the end-to-end success indicator.
    \item Collect at least 500 pipeline executions per pipeline-length bin, with bins at $n \in \{5, 10, 15, 20, 30, 50\}$.
    \item For each bin, estimate $\hat{p}$ (the empirical per-step success rate) and $\hat{R}(n)$ (the empirical end-to-end success rate).
    \item Additionally, collect data on step-to-step error correlation by recording whether failure at step $i$ predicts failure at step $i+1$ (beyond the base rate). This tests the independence assumption.
\end{enumerate}

\paragraph{Testable prediction.}
\begin{itemize}[nosep]
    \item Plot $\ln \hat{R}(n)$ against $n$. Under the model, this should be approximately linear with slope $\ln \hat{p}$.
    \item For each pipeline-length bin, compute the empirical elasticity $\hat{\mathcal{E}} = \Delta \ln \hat{R} / \Delta \ln \hat{p}$ (using within-bin variation in $p$ due to model differences or task difficulty). The prediction is $\hat{\mathcal{E}} \approx n$.
    \item Specifically: for $n = 20$ and $\hat{p} = 0.95$, the predicted reliability is $0.95^{20} = 0.36$, with 95\% CI bounds derivable from the binomial model.
\end{itemize}

\paragraph{Falsification criterion.} The result is falsified if \emph{any} of the following hold:
\begin{enumerate}[nosep]
    \item The relationship between $\ln \hat{R}(n)$ and $n$ is not approximately linear (e.g., $R^2 < 0.7$ for the linear fit), indicating that the series reliability model is structurally wrong (not merely noisy).
    \item The empirical elasticity $\hat{\mathcal{E}}$ differs from $n$ by more than a factor of 2 across at least 3 pipeline-length bins, indicating that the elasticity claim is quantitatively misleading.
    \item Step-to-step error correlations are so strong (correlation coefficient $> 0.5$) that the independence model provides a poor approximation, in which case the common-cause model from the Reliability pillar must be used as the baseline.
\end{enumerate}

\paragraph{Expected timeline and resources.} Requires partnership with a harness vendor or access to open-source harness telemetry. Estimated: 3--6 months for instrumentation and data collection; 1--2 months for analysis. Minimum viable dataset: 3,000 pipeline executions across 3+ pipeline-length bins.

\subsubsection{Result 2: Optimal Context Budget (Theorem~\ref{thm:context-budget})}

\paragraph{Formal claim.} Under the power-law degradation model $\delta(n) = (n/W_{\text{eff}})^{\alpha_d}$, the optimal context size satisfies $|C^*| < W$, with practical estimates $|C^*| \approx 0.15W$ to $0.40W$.

\paragraph{Data to collect.}
\begin{enumerate}[nosep]
    \item Select a benchmark of 200+ coding tasks spanning five difficulty levels and three code domains (web, systems, data). Tasks must have known-correct solutions for automated evaluation.
    \item For each task, prepare context sets at 10 sizes: 5\%, 10\%, 15\%, 20\%, 30\%, 40\%, 50\%, 60\%, 80\%, 100\% of the model's effective context window $W_{\text{eff}}$.
    \item Context content at each size level should be selected by the same retrieval algorithm (e.g., BM25 + reranking), adding lower-relevance items as size increases, to ensure that the size variation reflects the natural relevance gradient.
    \item For each (task, context-size) pair, run the agent 5 times (to estimate variance) and record: (a) task success (binary), (b) code quality score (0--10, via automated rubric), (c) token usage.
    \item Repeat across at least 3 frontier models from different providers.
\end{enumerate}

\paragraph{Testable prediction.}
\begin{itemize}[nosep]
    \item Plot mean task success against context size (as fraction of $W_{\text{eff}}$). Under the model, this curve should be unimodal: increasing for small context (more relevant information), peaking at $|C^*|/W_{\text{eff}} \in [0.15, 0.40]$, and declining for large context (degradation dominates).
    \item The peak location should be robust across models (within 15 percentage points), though the peak height will vary.
    \item For tasks where the model's prior is strong (boilerplate, common patterns), the peak should shift left (less context needed). For tasks requiring novel logic, the peak should shift right.
\end{itemize}

\paragraph{Falsification criterion.} The result is falsified if \emph{any} of the following hold:
\begin{enumerate}[nosep]
    \item Task success is monotonically non-decreasing in context size for a majority ($>60\%$) of tasks and models, indicating that degradation is not a binding constraint in practice.
    \item The peak, when it exists, occurs above $0.60 W_{\text{eff}}$ for the majority of tasks, indicating that the $[0.15, 0.40]$ range is too conservative.
    \item The peak location varies by more than 40 percentage points across models (e.g., 0.10 for one model and 0.55 for another), indicating that the ``optimal budget'' is model-specific rather than a general property of context selection.
\end{enumerate}

\paragraph{Expected timeline and resources.} Can be conducted with API access to frontier models. Estimated: 2--4 months for benchmark construction and execution; 1 month for analysis. Cost: approximately \$5,000--\$15,000 in API fees depending on model pricing and repetition count.

\subsubsection{Result 3: Structural Enforcement Threshold (Theorem~\ref{thm:structural-boundary})}

\paragraph{Formal claim.} Instructional enforcement (prompts, rules files) achieves compliance $(1-\epsilon)^T$, decaying exponentially in pipeline length $T$. Structural enforcement achieves compliance 1. The threshold for cost-effectiveness is $L > L^* = (c_s - c_p)/\epsilon$, which drops as $L^*/n$ for multi-step pipelines.

\paragraph{Data to collect.}
\begin{enumerate}[nosep]
    \item Define 10 invariants spanning the decidability spectrum: 4 decidable (e.g., ``never import package X,'' ``all functions have type annotations,'' ``no files exceed 500 lines,'' ``API keys never appear in source''); 3 semi-decidable (e.g., ``all error paths return appropriate status codes,'' ``no race conditions in shared state access,'' ``all user inputs are sanitized''); 3 undecidable (e.g., ``no unnecessary abstraction layers,'' ``naming conventions reflect domain semantics,'' ``code follows the existing architectural style'').
    \item For each invariant, implement both an instructional enforcement mechanism (system prompt instruction, CLAUDE.md rule, or rules file entry) and, where possible, a structural enforcement mechanism (linter rule, pre-commit hook, file-system permission, CI gate).
    \item Run a standardized coding agent through 50-step pipelines on a realistic codebase, with $T \in \{10, 20, 30, 50, 75, 100\}$ steps. At each step, record whether the invariant holds. Use 3 different frontier models.
    \item For each (invariant, enforcement-type, model, $T$) combination, compute the compliance rate. Minimum 30 trials per combination to achieve adequate statistical power.
\end{enumerate}

\paragraph{Testable prediction.}
\begin{itemize}[nosep]
    \item For structural enforcement on decidable invariants, compliance should be $\geq 0.99$ regardless of $T$ and model.
    \item For instructional enforcement, compliance should decay exponentially. Specifically, plotting $\ln(\text{compliance})$ against $T$ should yield an approximately linear relationship with slope $\approx \ln(1-\hat{\epsilon})$, from which $\hat{\epsilon}$ can be estimated.
    \item The estimated $\hat{\epsilon}$ should fall in the range $[0.005, 0.10]$ for decidable invariants and $[0.05, 0.30]$ for semi-decidable invariants.
    \item The compliance gap between structural and instructional enforcement should be statistically significant ($p < 0.01$) for $T \geq 20$.
\end{itemize}

\paragraph{Falsification criterion.} The result is falsified if \emph{any} of the following hold:
\begin{enumerate}[nosep]
    \item Instructional compliance does not decay with $T$ (i.e., compliance is approximately constant across pipeline lengths), indicating that the $(1-\epsilon)^T$ model is wrong and agents may have mechanisms (e.g., in-context learning, self-correction) that counteract exponential decay.
    \item Structural enforcement on decidable invariants fails at rates $> 5\%$ (due to agent workarounds, tool limitations, or enforcement gaps), indicating that the ``compliance = 1'' claim is practically false.
    \item The compliance gap between structural and instructional enforcement is not statistically significant for $T \leq 50$, indicating that the theoretical advantage of structural enforcement is too small to matter in practice.
\end{enumerate}

\paragraph{Expected timeline and resources.} Requires moderate engineering to implement enforcement mechanisms and instrumentation. Estimated: 4--6 months including implementation and data collection. Minimum viable dataset: 10 invariants $\times$ 2 enforcement types $\times$ 3 models $\times$ 6 pipeline lengths $\times$ 30 trials = 10,800 pipeline executions.

\subsubsection{Result 4: Optimal Agent Count (Theorem~\ref{thm:extended-amdahl})}

\paragraph{Formal claim.} For $n$ agents with per-agent defect injection rate $\delta_a$ and serial fraction $s$, effective speedup is $S_{\text{eff}}(n) = [1/(s + (1-s)/n)] \cdot (1-\delta_a)^n$, with optimal agent count $n^* \approx 1/\sqrt{\delta_a}$.

\paragraph{Data to collect.}
\begin{enumerate}[nosep]
    \item Select 30+ software engineering tasks of moderate complexity (2--8 hours of human developer time), covering feature implementation, bug fixing, and refactoring. Tasks must be decomposable into parallel subtasks of varying granularity.
    \item For each task, execute with $k \in \{1, 2, 3, 4, 5, 7, 10, 15, 20\}$ agents. Use at least two coordination topologies: (a) independent (no inter-agent communication) and (b) centrally coordinated (orchestrator assigns subtasks, reviews outputs, resolves conflicts).
    \item For each (task, $k$, topology) configuration, measure: (a) wall-clock time $T_k$, (b) code quality score $Q_k$ (via automated tests + human review on a subsample), (c) merge conflict count and resolution time, (d) rework iterations (code that was written, then reverted or rewritten).
    \item Compute Quality-Adjusted Speedup: $\text{QAS}(k) = (T_1 \cdot Q_1) / (T_k \cdot Q_k)$.
    \item Estimate $\hat{\delta}_a$ from the per-agent defect rate and $\hat{s}$ from the serial fraction of each task (via dependency analysis of the task decomposition).
\end{enumerate}

\paragraph{Testable prediction.}
\begin{itemize}[nosep]
    \item QAS should be unimodal in $k$: increasing for $k \leq n^*$, then decreasing.
    \item The peak should occur at $\hat{k}^* \approx 1/\sqrt{\hat{\delta}_a}$, within a factor of 2.
    \item For independent coordination with typical $\delta_a \approx 0.03$--$0.07$, the peak should be at $k^* \in [3, 6]$.
    \item At $k = 2n^*$, QAS should be measurably below its peak value (by at least 20\%), confirming that over-parallelization is costly.
    \item For the centrally coordinated topology, $\delta_a$ should be lower (better error correction), yielding a higher $n^*$.
\end{itemize}

\paragraph{Falsification criterion.} The result is falsified if \emph{any} of the following hold:
\begin{enumerate}[nosep]
    \item QAS is monotonically non-decreasing in $k$ up to $k = 20$ for a majority of tasks, indicating that the reliability penalty $(1-\delta_a)^n$ does not dominate within the tested range and the optimal agent count is much higher than predicted.
    \item The observed peak $\hat{k}^*$ differs from $1/\sqrt{\hat{\delta}_a}$ by more than a factor of 3 across a majority of tasks, indicating that the Extended Amdahl model is a poor predictor.
    \item QAS $< 1$ for $k = 2$ (i.e., even minimal parallelism is net-negative) for a majority of tasks and topologies, indicating a qualitative failure of the parallelism model that suggests coordination overhead dominates even before reliability penalties matter.
\end{enumerate}

\paragraph{Expected timeline and resources.} The most resource-intensive experiment in this roadmap, requiring either a multi-agent harness platform or substantial API costs. Estimated: 6--12 months. Cost: \$20,000--\$50,000 in compute. Minimum viable dataset: 30 tasks $\times$ 9 agent counts $\times$ 2 topologies = 540 configurations, each with quality evaluation.

\subsubsection{Result 5: Governance Capacity Bound (Theorem~\ref{thm:governance-capacity})}

\paragraph{Formal claim.} If the architectural drift rate $R_{\text{drift}}$ (bits per unit time) exceeds the governance channel capacity $C_{\text{gov}}$ (bits per unit time), no governance scheme prevents unbounded divergence from the intended architecture.

\paragraph{Data to collect.}
\begin{enumerate}[nosep]
    \item Select 20+ open-source repositories that (a) have Architecture Decision Records or equivalent governance artifacts, (b) have measurable AI adoption (via commit metadata, co-author tags, or tool integration logs), and (c) span at least 12 months of history pre- and post-AI adoption.
    \item For each repository, construct proxy measures for the two key quantities:
    \begin{itemize}[nosep]
        \item \textbf{Drift rate proxy} ($\hat{R}_{\text{drift}}$): rate of architectural violations per month, measured as: (a) ADR compliance violations (commits that contradict accepted ADRs, detected by automated pattern matching), (b) layer boundary violations (imports across declared layer boundaries), (c) dependency direction violations (detected by dependency analysis tools), (d) style/convention violations in new code. Each violation category receives a weight proportional to its severity. The composite is normalized to bits per month using $\hat{R}_{\text{drift}} = \sum_i w_i \cdot r_i$, where $r_i$ is the violation rate for category $i$.
        \item \textbf{Governance capacity proxy} ($\hat{C}_{\text{gov}}$): governance throughput per month, measured as: (a) code review throughput (PRs reviewed per month $\times$ mean review depth), (b) ADR update rate, (c) architectural review meeting frequency $\times$ mean decision output, (d) automated governance gate coverage (fraction of invariants enforced structurally). The composite is $\hat{C}_{\text{gov}} = \sum_j v_j \cdot g_j$, where $g_j$ is the governance activity rate for category $j$.
    \end{itemize}
    \item For each repository, measure \textbf{architectural health} over time using: (a) coupling density (mean module coupling, measured monthly), (b) abstraction stability (rate of interface changes per month), (c) technical debt proxies (TODO density, code duplication rate, cyclomatic complexity growth rate), (d) developer satisfaction with architecture (from surveys, if available).
    \item Compute the ratio $\hat{R}_{\text{drift}} / \hat{C}_{\text{gov}}$ monthly for each repository and plot against architectural health metrics.
\end{enumerate}

\paragraph{Testable prediction.}
\begin{itemize}[nosep]
    \item Repositories where $\hat{R}_{\text{drift}} / \hat{C}_{\text{gov}} > 1$ (drift exceeds governance) for sustained periods ($\geq$ 3 consecutive months) should show measurable architectural degradation: increasing coupling density, rising duplication, growing technical debt proxies.
    \item Repositories where $\hat{R}_{\text{drift}} / \hat{C}_{\text{gov}} < 1$ consistently should show stable or improving architectural health.
    \item The transition should be relatively sharp: repositories near the threshold ($0.8 < \hat{R}_{\text{drift}} / \hat{C}_{\text{gov}} < 1.2$) should show mixed signals, while those clearly above or below should show consistent degradation or stability.
    \item Repositories that increased AI adoption (higher commit velocity) without increasing governance capacity should show a transition from sub-threshold to super-threshold, with architectural degradation following the transition with a lag of 2--6 months.
\end{itemize}

\paragraph{Falsification criterion.} The result is falsified if \emph{any} of the following hold:
\begin{enumerate}[nosep]
    \item No measurable relationship exists between the $\hat{R}_{\text{drift}} / \hat{C}_{\text{gov}}$ ratio and architectural health metrics (correlation $< 0.2$ across repositories), indicating that the governance capacity framing has no predictive power.
    \item Repositories sustain $\hat{R}_{\text{drift}} / \hat{C}_{\text{gov}} > 1.5$ for 6+ months without measurable architectural degradation, indicating that the ``unbounded divergence'' prediction is too pessimistic and self-correcting mechanisms (developer workarounds, organic refactoring, competitive pressure) prevent divergence even without formal governance capacity.
    \item The proxy measures $\hat{R}_{\text{drift}}$ and $\hat{C}_{\text{gov}}$ cannot be operationalized with inter-rater reliability $> 0.6$ (Cohen's $\kappa$), indicating that the theoretical quantities are too vague to measure, which would itself be a significant finding about the framework's practical applicability.
\end{enumerate}

\paragraph{Expected timeline and resources.} The most methodologically challenging experiment, requiring careful proxy construction and validation. Estimated: 12--18 months for the longitudinal component; initial cross-sectional results possible in 6 months. Requires access to repository history and governance artifacts for 20+ projects. No significant compute cost, but substantial research analyst time.

\paragraph{Summary of the verification roadmap.} Table~\ref{tab:roadmap-summary} summarizes the five experiments.

\begin{table}[H]
\centering
\caption{Verification roadmap summary. Each row identifies the formal result, the core prediction, the falsification standard, and the estimated resource requirement.}
\label{tab:roadmap-summary}
\small
\begin{tabular}{p{2.5cm}p{3.5cm}p{4cm}p{2.5cm}}
\toprule
\textbf{Result} & \textbf{Core Prediction} & \textbf{Falsification Standard} & \textbf{Est. Resources} \\
\midrule
Compound error sensitivity & $\ln R$ linear in $n$; elasticity $= n$ & $R^2 < 0.7$ or elasticity off by $>2\times$ & 3--6 months; vendor partnership \\
\addlinespace
Optimal context budget & Unimodal success in context size; peak at $[0.15, 0.40]W$ & Monotonic success or peak $> 0.60W$ for majority & 2--4 months; \$5--15K API cost \\
\addlinespace
Structural enforcement threshold & Instructional compliance decays as $(1-\epsilon)^T$; structural = 1 & No decay with $T$, or structural fails $> 5\%$ & 4--6 months; moderate engineering \\
\addlinespace
Optimal agent count & QAS unimodal; peak at $\approx 1/\sqrt{\delta_a}$ & Monotonic QAS or peak off by $> 3\times$ & 6--12 months; \$20--50K compute \\
\addlinespace
Governance capacity & Drift/capacity ratio predicts arch. health & Correlation $< 0.2$ or sustained excess without degradation & 12--18 months; analyst time \\
\bottomrule
\end{tabular}
\end{table}

\subsection{Uncomputability of Idealized Metrics}

Several foundational quantities in the framework are uncomputable:

\begin{itemize}
    \item $K(P \mid S)$ (Kolmogorov complexity) is uncomputable by definition. The Shannon operationalization $H(P \mid S)$ is computable but distribution-dependent.
    \item The decidability classification (Section~\ref{sec:quality}) establishes that 4 of 18 slop types are undecidable in general. No automated system can achieve complete detection for these types.
    \item The Detection Ceiling (Theorem~\ref{thm:detection-ceiling}) is an information-theoretic bound, but computing $I(X; D_j)$ for a specific defect type and feature set requires knowing the joint distribution, which is not available in practice.
\end{itemize}

These uncomputability results are not limitations of the framework; they are features. They identify hard boundaries that no engineering effort can overcome, focusing design attention on what is achievable.

\subsection{The Theory Preservation Problem}

The Governance pillar identifies theory preservation as fundamentally open. If the Tacit Divergence Conjecture (Conjecture~\ref{conj:tacit}) holds, then perfect externalization of program theory is impossible in principle. Current approaches (ADRs, ARCHITECTURE.md, inline documentation) capture explicit knowledge but not tacit knowledge.

This problem becomes more acute with AI agents: the ``theory'' held by an agent exists only during a context window and is lost when the conversation ends. There is no persistent memory that preserves the agent's understanding of design rationale, rejected alternatives, or the causal model of change propagation. The theory preservation problem is arguably the most important open problem in the field.

\subsection{Scale-Dependent Governance ROI}

The governance capacity bound (Theorem~\ref{thm:governance-capacity}) establishes a threshold, but the return on governance investment is highly scale-dependent. For small projects with low change velocity, the drift rate $R_{\text{drift}}$ is naturally low, and minimal governance suffices. For large projects with high AI-driven change velocity, governance investment must scale, but the cost of governance also scales. The optimal governance investment as a function of project scale is not well characterized.

\subsection{The Measurement Problem (Goodhart's Law)}

Several metrics in the framework (VIH, QAS, detection rates) are susceptible to Goodhart's Law: when a measure becomes a target, it ceases to be a good measure. If harnesses are optimized for VIH, agents may learn to produce code that passes verification gates without being genuinely correct. If QAS becomes a target, teams may game quality metrics. The framework does not address this reflexivity problem, and it is unclear whether any formal framework can.


% ====================================================================
% APPENDIX A: CROSS-PILLAR THEOREM INDEX (EXPANDED)
% ====================================================================
\begin{appendices}

\section{Cross-Pillar Theorem Index}
\label{app:theorems}

Table~\ref{tab:theorem-index} provides a comprehensive index of the significant theorems, propositions, definitions, and conjectures across all seven pillar papers. The index is organized by pillar and includes all results that are referenced in cross-pillar arguments or that constitute a pillar's primary formal contributions.

\begin{table}[H]
\centering
\caption{Cross-pillar theorem index. Each entry summarizes a key formal result from the corresponding pillar paper. Entries marked with $\dagger$ are referenced in cross-pillar arguments in Section~\ref{sec:synthesis}.}
\label{tab:theorem-index}
\small
\begin{tabular}{p{1.8cm}p{3.8cm}p{7.5cm}}
\toprule
\textbf{Pillar} & \textbf{Result} & \textbf{Summary} \\
\midrule

%% ABSTRACTION
\multirow{6}{1.8cm}{Abstraction}
    & Abstraction Gap$^\dagger$ & $\mathcal{G}(S,P) = K(P \mid S)$: conditional Kolmogorov complexity between spec and code. \\
    & Gap Properties & Minimality ($\mathcal{G} \geq 0$), maximality ($\mathcal{G} \leq K(P)$), monotonicity (refining $S$ can only decrease $\mathcal{G}$), additivity (gap decomposes over independent components). \\
    & NID Universality & Normalised information distance $d(S,P)$ is a universal metric (Li-Vitanyi). \\
    & Spec-Code Convergence$^\dagger$ & For incompressible $P$, any $S$ with $K(P \mid S) \leq c$ satisfies $|S| \geq |P| - O(\log |P|)$. \\
    & Refinement Lattice & Specifications form a complete lattice under refinement; Galois connection $(\alpha, \gamma)$ between spec and program lattices. \\
    & Agent Decomposition & $\mathcal{G}(S,P) \leq \mathcal{G}(S, \hat{S}) + \mathcal{G}(\hat{S}, P) + O(\log n)$: intermediate specs reduce total gap up to a logarithmic overhead. \\

\midrule

%% INFORMATION
\multirow{8}{1.8cm}{Information}
    & Submodularity$^\dagger$ & Context relevance $f(C) = I(C; P \mid S, \theta)$ is submodular under mild conditions; greedy achieves $(1-1/e)$ approximation. \\
    & Non-Monotonicity & The degradation-adjusted objective $\tilde{f}(C) = f(C) \cdot (1-\delta(|C|))$ is non-monotone; requires randomized greedy with $1/2$-approximation. \\
    & Gap Decomposition$^\dagger$ & $H(P|S) = H(P|S,C,\theta) + I(P;C|S,\theta) + I(P;\theta|S)$, exact equality via Shannon chain rule. \\
    & Position Dominance & Context item position dominates item relevance in determining effective contribution; positional recall follows a U-shaped curve. \\
    & Context Budget$^\dagger$ & $|C^*| < W$; optimal context is sub-capacity, approximately $0.15W$ to $0.40W$. \\
    & Curation Beats Accumulation & Curated context (active tier) outperforms accumulated context of the same size. \\
    & Enforcement Dominance$^\dagger$ & Structural compliance $= 1$; instructional compliance $= (1-\epsilon)^T$, exponential decay in pipeline length. \\
    & Fresh Context Dominance & Agent with fresh context outperforms agent with stale context by $\sim$17\% (``Fresh Eyes'' effect). \\

\midrule

%% RELIABILITY
\multirow{8}{1.8cm}{Reliability}
    & Series Reliability & $R(n) = \prod p_i$; for homogeneous pipelines, $R(n) = p^n$. \\
    & Compound Sensitivity$^\dagger$ & Elasticity $\mathcal{E}(R,p) = n$; 1\% per-step improvement yields $n$\% system improvement. \\
    & Weakest-Link Priority & Improve the step with lowest $p_i$ first; marginal gain is maximized at the weakest link. \\
    & Correlated Failure & Under Marshall-Olkin common-cause model: $R_{\text{sys}} = (1-\gamma)\prod(1-\alpha_i)$, where $\gamma$ is shared failure rate. \\
    & Optimal Checkpoints$^\dagger$ & $k^* \approx \sqrt{n(1-p)c_0/c_v}$; Young-Daly analog. Three rounds capture $>96$\% of detectable errors. \\
    & Diminishing Returns & Marginal value of the $k$-th checkpoint: $\Delta M(k) = N_0 v(1-v)^{k-1}$, geometric decay. \\
    & Circular Bias$^\dagger$ & $\beta = \beta_A(1-\beta_{A'})$; self-verification ($A = A'$) yields zero bias reduction on shared blind-spot categories. \\
    & Structural Boundary$^\dagger$ & Structural enforcement cost-effective when $L > L^* = (c_s - c_p)/\epsilon$; threshold drops as $L^*/n$ for multi-step pipelines. \\

\midrule

%% COORDINATION
\multirow{10}{1.8cm}{Coordination}
    & Error Amplification$^\dagger$ & Under independence, $A_e = (1-(1-\epsilon)^k)/\epsilon \to k$; empirical is $17.2\times$ (correlation multiplier). \\
    & Saturation Crossover & Scaling benefits saturate at approximately $P^* \approx 0.45$ of the task completion rate. (Conjecture.) \\
    & Extended Amdahl$^\dagger$ & $S_{\text{eff}}(n) = [1/(s+(1-s)/n)] \cdot (1-\delta_a)^n$; optimal $n^* \approx 1/\sqrt{\delta_a}$. \\
    & NP-Hardness & Balanced $k$-way partition of the coupling graph is NP-hard. \\
    & Cheeger Inequality & Spectral gap bounds partition quality: $h(G) \leq \sqrt{2\lambda_2}$ (discrete Cheeger). \\
    & Cut-Conflict Bridge$^\dagger$ & $\mathbb{E}[\text{conflicts}] \leq \alpha \cdot W_{\text{cut}}(\pi)$; merge conflicts bounded by cut weight. \\
    & Contract Reduction & File-ownership contracts reduce conflict growth from $O(k^2)$ to $O(k)$. \\
    & Write Skew & Two agents with individually valid changes can produce an invalid combination; requires assume-guarantee contracts. \\
    & QAS$^\dagger$ & Quality-Adjusted Speedup $= (T_1 Q_1)/(T_k Q_k)$; QAS $< 1$ when quality degradation outweighs speed gain. \\
    & Conway's Law & Low-overhead coordination requires $G_{\text{coupling}} \subseteq G_{\text{comm}}$; graph containment condition. \\

\midrule

%% TEMPORAL
\multirow{8}{1.8cm}{Temporal}
    & Context Fidelity$^\dagger$ & $\Phi(t) = e^{-\Lambda(t-t_0)}$; exponential staleness decay with half-life $t_{1/2} = \ln 2 / \Lambda$. \\
    & VIH Decomposition$^\dagger$ & $\text{VIH} = \lambda_{\text{raw}} \cdot p_{\text{verify}}$; verified iterations per hour. \\
    & VIH Optimality & VIH maximized at unit elasticity: $\partial \ln \lambda / \partial v = -\partial \ln p_{\text{verify}} / \partial v$. \\
    & Speculation Test$^\dagger$ & Speculate iff $p/q > R/B$, where $R$ is rollback cost and $B$ is early-completion benefit. \\
    & Speculation Depth Limit & $k^* < \ln(1 + B_0/R_0)/\ln(1/p)$; depth beyond this has negative expected value. \\
    & Cache EOQ$^\dagger$ & $T^* = \sqrt{2R/(\lambda\alpha \cdot E_{\text{stale}})}$; optimal cache refresh interval. \\
    & Bayesian Mode Selection & Switch from fast to governed mode when $P(\text{Governed} \mid \mathbf{x}) > c_{GF}/(c_{FG} + c_{GF})$. \\
    & Pareto Convexification & Mixing fast and governed modes convexifies the speed-quality frontier; VIH optimum lies at the tangency of the VIH iso-hyperbola. \\

\midrule

%% QUALITY
\multirow{12}{1.8cm}{Quality}
    & Decidability Classes$^\dagger$ & 7 Decidable, 7 Semi-decidable, 4 Undecidable slop types across 18 AI code defect categories. \\
    & Latent Factor Structure & Three approximately orthogonal generative factors: Context Blindness, Completion Bias, Training Distribution Leakage. \\
    & Defense NP-Hardness$^\dagger$ & Layered defense optimization is NP-hard (via Weighted Set Cover); greedy achieves $(1-1/e)$ approximation. \\
    & Detection Ceiling$^\dagger$ & $d_{\ell,j} \leq I(X; D_j)/H(D_j)$ via Data Processing Inequality; information-theoretic limit on detection. \\
    & No Free Lunch & Context-poor layers cannot exceed the detection ceiling; context access is necessary for high detection rates. \\
    & FKG Correlation$^\dagger$ & $P(\text{both miss}) \geq P(\text{producer misses}) \cdot P(\text{judge misses})$; independence underestimates escape probability. \\
    & Diversity Gain & Weakly independent diverse verifiers outperform strongly correlated single-family verifiers. \\
    & ROI Ordering & Optimal layer ordering: decreasing $r_\ell/c_\ell$ (detection rate per unit cost). \\
    & Enforcement Gap & Syntactic properties enforceable with $P = 1$; semantic properties strictly less (Rice's theorem). \\
    & Compound Degradation$^\dagger$ & CI-passing changes yield superlinear entropy growth: $H \geq n\Delta H + \binom{n}{2}\Delta H_{\text{cross}}$. \\
    & Ratchet Effect & Codebase entropy is monotonically non-decreasing under CI-only governance. \\
    & Optimal Quality Level & Closed-form $q^* = (1/\beta)\ln(c_F \lambda_0 \beta / (c_P + c_A))$; higher for decidable types, lower for undecidable. \\

\midrule

%% GOVERNANCE
\multirow{11}{1.8cm}{Governance}
    & Theory Extraction Bound & Rate-distortion: $R(D) = \min_{p(\hat{T}|T): \mathbb{E}[d(T,\hat{T})] \leq D} I(T;\hat{T})$; minimum description length of program theory. \\
    & Tacit Divergence & Conjectured: $R(D) \to \infty$ as $D \to 0$ for tacit knowledge components; perfect externalization requires infinite description. \\
    & Incompressibility & Kolmogorov-random tacit components satisfy $K(t_\tau) \geq |t_\tau| - O(1)$; no lossless compression possible. \\
    & Double Bottleneck & $I(T;G) \leq I(T;(D,A,R)) \leq I(T;L)$; governance artifacts are twice-bottlenecked. \\
    & Capacity Bound$^\dagger$ & If $R_{\text{drift}} > C_{\text{gov}}$, no governance scheme prevents unbounded divergence. \\
    & AI Velocity Corollary & Increasing agent throughput $\lambda$ eventually overwhelms any fixed $C_{\text{gov}}$; governance must co-scale with velocity. \\
    & Cascade Stability & Three-loop cascade governance is stable iff each loop satisfies Nyquist-like sampling requirement: $f_k \geq 2\omega_d^{(k)}$. \\
    & Governance Bifurcation & Transcritical bifurcation at $\mu G / \gamma R = 1$; below: self-correcting; above: unbounded drift. \\
    & Ratchet Convergence$^\dagger$ & Closure operator on constraint lattice converges in $\leq |\mathcal{R}|$ steps. \\
    & Ossification & When contradictory constraints accumulate, the allowed change set becomes empty; the ratchet becomes pathological. \\
    & Decision Survival & Competing-risks Weibull model; median lifetimes: technology 2--4y, framework 3--5y, architecture 5--10y. \\

\bottomrule
\end{tabular}
\end{table}

The index contains 63 entries across 7 pillars. Of these, 19 (marked with $\dagger$) participate directly in the cross-pillar arguments developed in Section~\ref{sec:synthesis}. The remaining 44 are pillar-internal results that provide the technical foundations for the cross-pillar theorems. The density of cross-references (19 out of 63 results, approximately 30\%) confirms that the pillars are not independent theories but a connected formal network.

\end{appendices}
